%%%%%%%%%%%%%%%%%%%%%%%%%%%%%%%%%%%%%%%%%%%%%%%%%%%%%%%%%%%%%%%%%%%%%%
%  第 5 个逻辑月 · 月度报告（待做）
%  主题：整板验证
%
%  说明：组长将此月标注为「待做」。本报告以「验证方案设计」形式组织，
%        说明"验证什么、怎么验证、需要什么条件"，而非汇报已完成的结果。
%  编译：xelatex 本文件（跑两遍）
%%%%%%%%%%%%%%%%%%%%%%%%%%%%%%%%%%%%%%%%%%%%%%%%%%%%%%%%%%%%%%%%%%%%%%
\documentclass[stage2,fontset=fandol]{ydzjreport}

\stagetwosetup
  {第 5 个逻辑月 · 整板验证（待做）}%
  {冯旭东、陈思文、张昊洋}%
  {复合微通道分区冷却}%
  {【待确认】逻辑月 5 对应的自然月}%
  {整板集成；分区方案验证；温度均匀性；分流校核}%
  {本逻辑月为待做阶段。本报告给出整板验证的方案设计——
   验证什么、用什么指标、需要哪些前置条件。}%

\begin{document}
\makestagetwohead

\begin{notebox}
\textbf{本月状态}：组长提纲中本月标注为\textbf{「待做」}。
因此本报告不是"已完成工作的汇报"，而是\textbf{验证方案的设计}：
把"整板验证"这个较宽的说法，拆解为可检查的具体项目。
【待确认7】验证指标、仿真还是实物、三人分工。
\end{notebox}

% ==================================================================
\section{本月研究任务}

\subsection{任务目标}

把"整板验证"拆解为五项可检查的验证内容：

\begin{enumerate}
  \item \textbf{分区方案是否有效}——整板温度均匀性与热点是否受控；
  \item \textbf{流动分配是否合理}——各支路流量偏差与连接段压降；
  \item \textbf{能效是否可接受}——整板泵功、综合指标或热阻；
  \item \textbf{是否优于单流道}——与上逻辑月定型的单元结果对比；
  \item \textbf{控制是否可行}——执行器与散热片的特性辨识能否落地。
\end{enumerate}

\subsection{完成情况}

\begin{table}[H]
\centering
\caption{本月任务完成情况}
\label{tab:tasks}
\begin{tabularx}{\textwidth}{@{}Y c Y@{}}
\toprule
\rowcolor{ydzjlight} 验证内容 & 状态 & 说明 \\
\midrule
分区方案有效性 & 未开始 & 依赖网络与连接段方案 \\
流动分配合理性 & 未开始 & 依赖支路流量数据 \\
整板能效 & 未开始 & 依赖整板组合模型 \\
与单流道对比 & 未开始 & 单元基准已有，整板数据待出 \\
控制可行性 & 未开始 & 依赖执行器特性辨识实验 \\
\bottomrule
\end{tabularx}
\end{table}

\subsection{前置条件}

整板验证依赖前面几个逻辑月的产出，前置条件如下：

\begin{table}[H]
\centering
\caption{整板验证的前置条件}
\label{tab:pre}
\begin{tabularx}{\textwidth}{@{}Y c Y@{}}
\toprule
\rowcolor{ydzjlight} 前置条件 & 状态 & 说明 \\
\midrule
小块划分定义 & 【待确认9】 & 决定整板分区与网络形式 \\
单元级性能基准 & 已有 & 四种流道 39 工况、六种翼型 \\
网格基准 & 已有 & 约 150 万网格 \\
可信物性 UDF & 已编译，待挂接 & 需完成复算 \\
网络拓扑与连接段 & 【待补】 & 上逻辑月未完成 \\
支路流量数据 & 【待补】 & 依赖连接段建模 \\
执行器特性 & 【缺数据】 & 需实测辨识 \\
\bottomrule
\end{tabularx}
\end{table}

% ==================================================================
\section{验证方案设计}

\subsection{验证一：分区方案是否有效}

\textbf{要回答的问题}：采用差异化结构的分区冷板，
是否真的改善了温度分布，而不是仅仅改变了平均值。

\textbf{指标}：
\begin{itemize}
  \item 整板\textbf{最高温度}——是否低于允许上限；
  \item 整板\textbf{温度均匀性}——用同一指定表面的
        最高温与最低温之差，或温度标准差表征；
  \item 各区\textbf{实际温度}与设计预期的偏差。
\end{itemize}

\textbf{判据}：与均匀布置方案对比，在相同泵功下温度均匀性应明显改善。

\subsection{验证二：流动分配是否合理}

\textbf{要回答的问题}：并联或混合网络下，各小块是否都得到了设计流量。

\textbf{指标}：
\begin{itemize}
  \item 各\textbf{支路流量}相对设计值的偏差；
  \item \textbf{连接段压降}及其占整板压降的比例；
  \item 各小块\textbf{入口速度分布}的均匀性。
\end{itemize}

\textbf{判据}：无支路出现明显供液不足；入口速度分布无严重偏流。

\subsection{验证三：能效是否可接受}

\textbf{指标}：整板泵功、综合性能指标、峰值热阻。

\textbf{注意}：综合指标必须与单元级使用\textbf{同一基准}，
否则整板与单元不可比。

\subsection{验证四：是否优于单流道}

\textbf{要回答的问题}：分区复合方案相对单一流道，收益是否成立。

\textbf{指标}：相同热载荷下的最高温度、温度均匀性、泵功三项对比。

\textbf{注意}：上逻辑月已确认\textbf{单元排序不能直接搬到整板}，
因此本项对比必须用整板数据，不能引用单元级结论。

\subsection{验证五：控制是否可行}

\textbf{要回答的问题}：控制方向的两个待辨识函数能否落地。

\textbf{要辨识的对象}：
\begin{itemize}
  \item 执行器到流速的关系（泵转速、阀开度 → 流速）；
  \item 流速与水温到芯片温度的关系。
\end{itemize}

\textbf{所需条件}：泵的型号与流量范围、是否已加装水阀、
有无流量计、温度采样频率，
以及\textbf{各小块之间是否独立供液}（决定控制架构）。

【缺数据：上述控制实验条件】
【待确认8】控制方向所需数据的提供方式与格式。

% ==================================================================
\section{问题与对策}

\subsection{遇到的主要问题}

\begin{enumerate}
  \item \textbf{前置条件未齐。} 网络拓扑与连接段方案尚未完成，
        整板模型无法组装。
  \item \textbf{物性尚未挂接。} 整板验证若继续使用旧物性，
        结论的绝对数值不可信。
  \item \textbf{控制实验条件不明。} 泵、阀、流量计与供液方式
        均需确认后才能设计辨识实验。
\end{enumerate}

\subsection{下一步对策}

\begin{enumerate}
  \item 按上逻辑月计划先完成网络拓扑与连接段，再组装整板模型；
  \item 在整板计算前完成新 UDF 挂接与校验；
  \item 向组长与结构方向确认控制实验条件，
        特别是\textbf{各小块是否独立供液}。
\end{enumerate}

% ==================================================================
\section{阶段性结论}

\subsection{可以确定的}

\begin{enumerate}
  \item "整板验证"可拆解为五项可检查的内容，每项都有明确的指标与判据。
  \item 整板验证的\textbf{关键路径}是：小块划分定义 → 网络拓扑 →
        连接段 → 支路流量 → 整板验证。
  \item 单元排序不能直接搬到整板，因此整板验证\textbf{必须独立进行}。
\end{enumerate}

\subsection{尚不能确定的}

\begin{enumerate}
  \item \textbf{验证用仿真还是实物}，以及三人在其中的分工。
  \item \textbf{整板温度均匀性的合格判据}（多均匀算合格）。
  \item \textbf{控制实验的台架条件}，特别是小块是否独立供液。
\end{enumerate}

% ==================================================================
\section{下月计划}

\begin{table}[H]
\centering
\caption{下月工作计划}
\label{tab:plan}
\begin{tabularx}{\textwidth}{@{}Y Y c@{}}
\toprule
\rowcolor{ydzjlight} 事项 & 预期产出 & 时间 \\
\midrule
确定小块划分定义与网络拓扑 & 划分定义与拓扑方案 & 第 1 周 \\
组装整板组合模型并挂接新物性 UDF & 整板模型 & 第 1--2 周 \\
整板分流校核与温度均匀性计算 & 分流与温度数据 & 第 2--3 周 \\
与单流道方案的对比分析 & 对比结论 & 第 3--4 周 \\
控制实验条件确认与辨识方案 & 辨识实验方案 & 第 4 周 \\
\bottomrule
\end{tabularx}
\end{table}

% ==================================================================
\begin{thebibliography}{9}
\bibitem{proposal} 大创项目申报书：研究内容与实施方案. 2026.
\bibitem{data} 四种流道数据与折线图. 项目计算资料, 2026.
\bibitem{udf} 项目组. water\_properties\_refprop\_fit.c：液态水物性拟合与 Fluent UDF 源码. 项目计算资料, 2026.
\end{thebibliography}

\end{document}
